% ================================================================
%  conteudo/conclusao.tex
% ================================================================

\section{Conclusão}

A prática realizada permitiu o isolamento de microrganismos a partir de amostras de uvas, com destaque para o crescimento de colônias com morfologia compatível com leveduras. O protocolo empregado, baseado na coleta por swab e na técnica de esgotamento por estrias em meio sólido, mostrou-se adequado para a obtenção de culturas puras a partir de amostras mistas de origem vegetal.

Os resultados obtidos evidenciaram a riqueza da microbiota nativa presente na superfície e na polpa das uvas, reforçando o potencial desse substrato como fonte de linhagens microbianas com interesse biotecnológico. A incubação a 30~°C por seis dias mostrou-se suficiente para o desenvolvimento das colônias, possibilitando a análise macroscópica das características morfológicas dos isolados.

Conclui-se que o experimento atingiu o objetivo proposto, contribuindo para a compreensão prática das técnicas fundamentais de isolamento microbiano abordadas na disciplina. Na sequência é possível até caracterizar os isolados obtidos, bem como a avaliação de suas propriedades fisiológicas e do potencial de aplicação em processos fermentativos industriais.
